% !TeX root = ../main.tex
% !TeX encoding = utf8
\section{ỨNG DỤNG BOOTSTRAP CI TRONG Y SINH}

\subsection{Bootstrap trong phân tích dữ liệu sống còn}

Dữ liệu sống còn có thể bị kiểm duyệt nếu thời gian theo dõi kết thúc trước khi biến cố xuất hiện hoặc người tham gia rời nghiên cứu. Trong mô hình Cox, Hazard Ratio (HR) là hàm mũ của hệ số hồi quy. Khi Bootstrap được dùng để đánh giá độ bất định của HR, cách lấy mẫu lại cần phù hợp với cấu trúc dữ liệu \cite{davison1997}.

Với hệ số ước lượng $\widehat{\beta}$, $\widehat{HR}=\exp(\widehat{\beta})$.
Khoảng Wald được tính trên thang log-HR rồi chuyển về thang HR:
\[
\left[
\exp(\widehat{\beta}-z_{1-\alpha/2}SE(\widehat{\beta})),
\exp(\widehat{\beta}+z_{1-\alpha/2}SE(\widehat{\beta}))
\right].
\]
Do phép biến đổi hàm mũ, hai cận của khoảng Wald Cox trên thang HR đều dương.
Normal Bootstrap CI tính trực tiếp trên thang HR là một cách xây dựng khác.

\subsection{Ví dụ giả định: ước lượng Hazard Ratio}

\textit{Ví dụ dưới đây dùng số liệu giả định để minh họa phép tính, không mô tả
một nghiên cứu hay thử nghiệm lâm sàng cụ thể.}

Quy trình minh họa gồm các bước sau:
\begin{enumerate}
	\item \textbf{Ước lượng Hazard Ratio:} Đặt $HR=h_{\text{điều trị}}/h_{\text{đối chứng}}$ và giả sử ước lượng trên mẫu gốc là $\widehat{HR}_n=1.85$. Nếu biến cố là tử vong, giá trị này biểu thị hazard tử vong tức thời ở nhóm điều trị cao hơn nhóm đối chứng theo mô hình Cox; cách diễn giải phụ thuộc vào quy ước mã hóa nhóm và biến cố.
	\item \textbf{Bootstrap phân tầng ($B=2000$):} Rút $2000$ mẫu Bootstrap phân tầng (stratified bootstrap), lấy mẫu riêng trong từng nhóm để giữ nguyên cỡ nhóm \cite{davison1997}. Với mỗi mẫu $b$, tính lại $\widehat{\text{HR}}^{*(b)}$.
	\item \textbf{Phân phối Bootstrap giả định:} Giả sử phân phối các bản sao $\widehat{\text{HR}}^*$ lệch phải và có các đặc trưng sau:
	\begin{itemize}
		\item Sai số chuẩn Bootstrap: $\widehat{se}^* = 0.42$.
		\item Phân vị dưới $2.5\%$: $q_{0.025}^* = 1.15$.
		\item Phân vị trên $97.5\%$: $q_{0.975}^* = 2.82$.
	\end{itemize}
	\item \textbf{Xây dựng và So sánh các Khoảng tin cậy $95\%$:}
	\begin{itemize}
		\item \textbf{Normal Bootstrap CI:}
		\begin{equation}
			CI_{\text{Normal}} = [1.85 - 1.96 \times 0.42, \quad 1.85 + 1.96 \times 0.42] = [1.027, \quad 2.673]
		\end{equation}
		\item \textbf{Basic Bootstrap CI (Quantile Flipping):}
		\begin{equation}
			CI_{\text{Basic}} = [2(1.85) - 2.82, \quad 2(1.85) - 1.15] = [0.880, \quad 2.550]
		\end{equation}
	\end{itemize}
\end{enumerate}

\textbf{Diễn giải:}
\begin{itemize}
	\item Hai khoảng thu được là
	$CI_{\text{Normal}}=[1.027,2.673]$ và
	$CI_{\text{Basic}}=[0.880,2.550]$. Trong ví dụ này, khoảng Normal
	không chứa $1$, còn khoảng Basic có chứa $1$.
	
	\item Hai khoảng đưa ra kết luận khác nhau về việc giá trị tham chiếu $1$
	có nằm trong khoảng hay không. Do số liệu được giả định, ví dụ này chỉ
	minh họa phép tính, không đánh giá độ chính xác hay xác suất bao phủ của
	hai phương pháp.
	
	\item Với dữ liệu sống còn thực tế, đơn vị lấy mẫu Bootstrap là bệnh nhân:
	mỗi lần lấy mẫu cần giữ đi cùng thời gian theo dõi, trạng thái biến cố
	và nhóm điều trị của bệnh nhân. Sau đó khớp lại mô hình Cox trên từng
	mẫu Bootstrap; cần báo cáo số mẫu không hội tụ nếu có.
\end{itemize}
\subsection{Bài báo khoa học: khoảng Bootstrap cho relative validity}

Deng và cộng sự (2013) dùng Bootstrap để đánh giá ý nghĩa thống kê của
relative validity (RV) và khảo sát các yếu tố ảnh hưởng đến kết luận.
Phân tích dựa trên câu trả lời của 453 bệnh nhân bệnh thận mạn đối với 16
thang đo đặc hiệu và tổng quát; RV cho biết mỗi thang đo phân biệt các
nhóm lâm sàng tốt đến mức nào so với thang đo tham chiếu \cite{deng2013}.

Đại lượng chính là relative validity (RV), được tính bằng tỷ số giữa
thống kê $F$ của công cụ so sánh và thống kê $F$ của công cụ tham chiếu:
\[
RV =
\frac{F_{\text{công cụ so sánh}}}
{F_{\text{công cụ tham chiếu}}}.
\]
Công cụ QDIS-CKD CAT-5 được chọn làm tham chiếu, có $F=57.43$ và
$RV=1$. Giá trị $RV<1$ biểu thị công cụ so sánh phân biệt các nhóm kém
hơn công cụ tham chiếu theo tiêu chí này.

Bài báo sử dụng khoảng tin cậy Bootstrap $95\%$ cho RV. Chú thích bảng
kết quả xác định đây là khoảng BCa (Bias-Corrected and Accelerated),
không phải Normal CI hay Basic CI. Bảng dưới đây trích một số kết quả
đại diện từ nghiên cứu \cite{deng2013}.

\begin{table}[htbp]
	\centering
	\small
	\renewcommand{\arraystretch}{1.2}
	\caption{Một số kết quả RV và khoảng tin cậy BCa trong nghiên cứu của Deng và cộng sự}
	\label{tab:deng-bca}
	\begin{tabularx}{\textwidth}{l c c c X}
		\toprule
		\textbf{Công cụ so sánh} &
		\textbf{$F$} &
		\textbf{RV} &
		\textbf{BCa 95\% CI} &
		\textbf{Diễn giải so với tham chiếu} \\
		\midrule
		QDIS-CKD Static-6 & 50.15 & 0.87 & [0.72; 1.03] &
		Khoảng chứa 1; chưa cho thấy khác biệt rõ ở mức 5\%. \\
		
		QDIS-CKD Static-34 & 48.01 & 0.84 & [0.71; 0.97] &
		Khoảng nằm dưới 1; RV thấp hơn tham chiếu theo khoảng BCa. \\
		
		KDQOL Burden & 44.46 & 0.77 & [0.53; 1.09] &
		Khoảng chứa 1; chưa cho thấy khác biệt rõ ở mức 5\%. \\
		
		KDQOL Symptoms & 15.11 & 0.26 & [0.13; 0.44] &
		Khoảng nằm dưới 1; RV thấp hơn tham chiếu theo khoảng BCa. \\
		\bottomrule
	\end{tabularx}
\end{table}

Bảng \ref{tab:deng-bca} cho thấy chỉ nhìn vào RV điểm là chưa đủ để
kết luận hai công cụ khác nhau. Chẳng hạn, RV của KDQOL Burden bằng
$0.77$, nhưng khoảng BCa $[0.53;1.09]$ vẫn chứa giá trị tham chiếu $1$.
Ngược lại, khoảng của QDIS-CKD Static-34 là $[0.71;0.97]$, không chứa
$1$. Vì vậy, theo cách diễn giải khoảng tin cậy trong nghiên cứu,
Static-34 có RV thấp hơn có ý nghĩa thống kê so với công cụ tham chiếu,
còn kết quả của KDQOL Burden chưa đủ để đưa ra kết luận tương tự
\cite{deng2013}.

Nghiên cứu còn thực hiện mô phỏng để khảo sát ảnh hưởng của cỡ mẫu,
thống kê $F$ của công cụ tham chiếu, tương quan giữa hai công cụ và số
lần lấy mẫu Bootstrap. Một kết quả đáng chú ý là khi $F$ của công cụ
tham chiếu khoảng $57$ và tương quan giữa hai công cụ bằng $0.7$, công
suất đạt khoảng $80\%$ để phát hiện RV xấp xỉ $0.6$. Tác giả cũng nhận
xét rằng tăng số lần lặp Bootstrap vượt quá $500$ đem lại ít thay đổi
trong các kịch bản mô phỏng của bài báo \cite{deng2013}.

Ví dụ này minh họa BCa được dùng để lượng hóa độ bất định của một tỷ số
thống kê phức tạp. Tuy nhiên, bài báo không phải là phép so sánh trực
tiếp BCa với Normal CI và Basic CI. Do đó, kết quả này bổ sung ví dụ ứng
dụng thực tế cho báo cáo, còn việc so sánh Normal CI và Basic CI vẫn
được thực hiện riêng trong phần mô phỏng của nhóm.

\begin{lstlisting}[style=Rstyle, caption={Vẽ lại các khoảng BCa được công bố trong bài báo}, label={lst:deng-bca}]
	measure <- c(
	"KDQOL Symptoms",
	"KDQOL Burden",
	"QDIS-CKD Static-34",
	"QDIS-CKD Static-6"
	)
	
	rv <- c(0.26, 0.77, 0.84, 0.87)
	lower <- c(0.13, 0.53, 0.71, 0.72)
	upper <- c(0.44, 1.09, 0.97, 1.03)
	y <- seq_along(rv)
	
	png("deng_bca_ci.png", width = 1800, height = 1100, res = 200)
	
	par(mar = c(5, 10, 3, 2))
	plot(
	rv, y,
	xlim = c(0, 1.2),
	ylim = c(0.5, 4.5),
	yaxt = "n",
	ylab = "",
	xlab = "Relative validity (RV) va khoang tin cay BCa 95%",
	main = "Khoang BCa cua mot so cong cu trong nghien cuu Deng et al. (2013)",
	pch = 19
	)
	
	axis(2, at = y, labels = measure, las = 1)
	segments(lower, y, upper, y, lwd = 2)
	segments(lower, y - 0.08, lower, y + 0.08, lwd = 2)
	segments(upper, y - 0.08, upper, y + 0.08, lwd = 2)
	abline(v = 1, lty = 2, col = "red", lwd = 2)
	
	legend(
	"bottomright",
	legend = "RV = 1 (muc tham chieu)",
	col = "red",
	lty = 2,
	lwd = 2,
	bty = "n"
	)
	
	dev.off()
\end{lstlisting}

\subsection{Mã R minh họa Bootstrap cho Hazard Ratio}

Mã ở Listing~\ref{lst:deng-bca} chỉ vẽ lại các ước lượng RV và cận BCa đã công bố; mã không tái lập phân tích từ dữ liệu gốc. Mã tiếp theo tạo dữ liệu giả định, khớp mô hình Cox và lấy mẫu lại phân tầng theo nhóm để minh họa việc tính Normal CI và Basic CI; đây không phải phân tích dữ liệu của nghiên cứu Deng và cộng sự:

\begin{lstlisting}[style=Rstyle, caption={Bootstrap HR}]
	# 1. Sinh du lieu gia dinh Thoi gian song sot va Trang thai (1: Bien co, 0: Khuyet)
	set.seed(2026)
	n <- 40
	group <- rep(c(0, 1), each = n/2) # 0: Control, 1: Treatment
	# Sinh thoi gian theo phan phoi Mu phu thuoc vao nhom
	time <- rexp(n, rate = 0.1 * exp(-0.6 * group))
	status <- rbinom(n, size = 1, prob = 0.85) # 85% quan sat duoc bien co
	
	data_med <- data.frame(time = time, status = status, group = group)
	
	# 2. Tinh Hazard Ratio (HR) tren mau goc qua mo hinh Cox
	library(survival)
	fit_orig <- coxph(Surv(time, status) ~ group, data = data_med)
	hr_hat <- exp(coef(fit_orig))
	
	# 3. Stratified Bootstrap Monte Carlo (B = 2000)
	B <- 2000
	boot_hrs <- numeric(B)
	
	for (b in 1:B) {
		# Rut mau phan tang theo nhom
		idx0 <- sample(which(data_med$group == 0), replace = TRUE)
		idx1 <- sample(which(data_med$group == 1), replace = TRUE)
		boot_data <- data_med[c(idx0, idx1), ]
		
		fit_boot <- coxph(Surv(time, status) ~ group, data = boot_data)
		boot_hrs[b] <- exp(coef(fit_boot))
	}
	
	# 4. Tinh dac trung va Khoang tin cay
	se_star <- sd(boot_hrs)
	q_low   <- quantile(boot_hrs, 0.025)
	q_upp   <- quantile(boot_hrs, 0.975)
	
	ci_normal <- c(hr_hat - 1.96 * se_star, hr_hat + 1.96 * se_star)
	ci_basic  <- c(2 * hr_hat - q_upp, 2 * hr_hat - q_low)
	
	cat("Original HR Hat:", round(hr_hat, 3), "\n")
	cat("Normal CI (95%): [", round(ci_normal[1], 3), ",", round(ci_normal[2], 3), "]\n")
	cat("Basic CI  (95%): [", round(ci_basic[1], 3), ",", round(ci_basic[2], 3), "]\n")
	
	# 5. Ve do thi phan phoi Bootstrap HR va Khoang tin cay
	hist(boot_hrs, breaks = 35, probability = TRUE, col = "aliceblue",
	main = "Phan phoi Bootstrap cho Hazard Ratio (HR) trong Nghien cuu Y sinh",
	xlab = "Gia tri Hazard Ratio Bootstrap (HR*)", border = "white")
	lines(density(boot_hrs), col = "darkblue", lwd = 2)
	
	# Ve duong nguong HR = 1.0 va vi tri HR Hat
	abline(v = 1.0, col = "gray40", lwd = 2, lty = 3) # Nguong khong co tac dung
	abline(v = hr_hat, col = "black", lwd = 2, lty = 1)
	
	# Ve thanh Khoang tin cay Normal CI vs Basic CI
	y_m <- max(density(boot_hrs)$y)
	lines(ci_normal, c(y_m * 0.75, y_m * 0.75), col = "red", lwd = 4)
	lines(ci_basic,  c(y_m * 0.55, y_m * 0.55), col = "green3", lwd = 4)
	
	legend("topright", 
	legend = c("HR Hat", "Nguong HR=1.0", "Normal CI", "Basic CI"),
	col = c("black", "gray40", "red", "green3"), lty = c(1, 3, 1, 1), lwd = c(2, 2, 4, 4))
\end{lstlisting}

% ==================== PHẦN 8 ====================
